% ─────────────────────────────────────────────────────────────
% Abstract
% ─────────────────────────────────────────────────────────────
\thispagestyle{plain}
\begin{center}
    {\fontsize{18pt}{22pt}\selectfont \textbf{Abstract} \par}
\end{center}

\vspace{0.5cm}
\noindent Active Directory Certificate Services (ADCS) forms the identity and cryptographic backbone of modern Windows enterprise networks, issuing public-key certificates used for domain authentication, transport encryption, and single sign-on. When certificate templates, issuance policies, and access control lists (ACLs) are misconfigured, unprivileged accounts can obtain unauthorized administrative certificates and take over entire Active Directory forests. While signature-based scanners (such as Certipy and BloodHound) identify known vulnerable flag combinations, they either evaluate configuration attributes in isolation or incur high traversal costs on dense enterprise graphs.

This thesis examines whether heterogeneous graph attention networks can reliably detect ADCS privilege escalation vectors (ESC1, ESC2, ESC3, ESC4, ESC9, and ESC13) directly from relational identity multigraphs. Using the experimental audit framework of Arp et al. (2022), we first evaluate synthetic Active Directory datasets and identify three systemic methodological pitfalls: positional identifier leakage in generation routines, feature asymmetry between baselines, and benchmark saturation where synthetic rules can be recovered by shallow tabular heuristics. We prove mathematically (Theorem 1) that residual skip-connections are indispensable for preventing representation collapse on identity graphs; without them, template representations become independent of their initial configuration flags, dropping classification F1 from 0.9986 to 0.4768.

On an audited benchmark of 700 enterprise environments evaluated under 5-fold cross-validation with Nadeau-Bengio corrected variance estimators, our model achieves a Macro-F1 of 0.9986, significantly outperforming flat classifiers (0.8600) and signature heuristics (0.7791). However, when tested against an adversarial Zero-Shot Hard Negative suite ($n=63$) containing non-exploitable templates with dangerous flags, pure neural models collapse to 1.59\% accuracy due to shortcut learning on local template flags. Symbolic path traversal, by contrast, correctly recognizes path absence in 84.13\% of cases.

To reconcile this gap, we develop a Two-Tier Neuro-Symbolic architecture. The neural tier acts as a high-throughput candidate screener, filtering out over 89\% of benign templates, while a localized symbolic oracle provides deterministic reachability proofs for flagged candidates. We show that false-positive suppression on hard negatives is achieved by design through this division of labor. Finally, we model identity privilege revocation as a Stackelberg security game and establish that optimal edge interdiction is NP-hard via reduction from Directed Multiway Cut, introducing a polynomial-time greedy heuristic that severs 92\% of reachable attack paths under bounded operational budgets.

\vspace{0.5cm}
\noindent\textbf{Keywords:} Active Directory Certificate Services (ADCS), Graph Attention Networks, Heterogeneous Graphs, Identity and Access Management, Neuro-Symbolic Security.
